\documentclass{ximera}
\title{Relative Motion, CW}

\newcommand{\pskip}{\vskip 0.1 in}

\begin{document}
\begin{abstract}
An introduction to relative motion.
\end{abstract}
\maketitle

\section{Relative Position}

\begin{example} \label{Ex:dlfkjef3}
For the following questions draw pictures with the non-colinear points $P$, $Q$, and $R$ along with the relevant vectors to help with your explanations.

\begin{enumerate}
\item What vector describes the position of $P$ relative to $Q$ and the position of $Q$ relative to $R$?

\item What vector describes the position of $R$ relative to $Q$?

\item Express the third vector above in terms of the first two.

\item Express the vector giving the position of $R$ relative to $Q$ in terms of the vectors giving the positions of $R$ relative to $P$ and of $Q$ relative to $P$. 
\end{enumerate}
\end{example}


\begin{example} \label{Ex:lferf303r0}
In this example, as in all others, we orient our Cartesian coordinate system with the positive $x$-axis pointing due east and the positive $y$-axis due north. We measure bearing as we do the polar angle, from the positive $x$-axis (due east) counterclockwise (toward the north) . Note that bearing is actually measured from the north toward the east. So a bearing of $\pi/2$ radians for us points due north, but in common use points due east.

Point $A$ is $6$ miles due east of point $B$. 

Point $C$ is $10$ miles from $A$ at a bearing of 2.4 radians.

\begin{enumerate}
\item Draw a reasonably accurate picture on a coordinate system as described above. Include the points $A$, $B$, $C$ and draw the two vectors giving the positions of $A$ relative to $B$ and of $C$ relative to $B$.

\item Draw a third vector giving the position of $C$ relative to $B$.

\item Express this third vector in terms of the first two without referencing the given information.

\item Find the exact components of the first two vectors. No calculator. Include units.

\item Use part (d) to find the exact components of the third vector. No calculator.

\item Find a simplified expression for the exact distance between $B$ and $C$. No calculator. Then use a calculator to approximate the distance.

\end{enumerate}
\end{example}


\section{Relative Motion}

\begin{example} \label{Ecr4322}
A planet with radius $a$ km revolves about its sun in a circular orbit of radius $b$ km at a constant speed. The equator of the planet lies in its orbital plane. There are four days in the planet's year and the planet rotates about its axis in the same sense as it revolves about its sun as illustrated below (where we see a cross-section of the sun and planet in the orbital plane).
\end{example}


Access Desmos interactive at
 
\href{https://www.desmos.com/calculator/golluueo9z}{Rotation Rates of Planets}

 
\begin{onlineOnly}
    \begin{center}
\desmos{golluueo9z}{900}{600}
\end{center}
\end{onlineOnly}


Suppose the point $Q$, at rest relative to the planet's equator, has the position shown at the start of the calendar year (time $t=0$).
\begin{enumerate}
\item Express the components of the vector
\[
  \overrightarrow{OP}  = \langle x , \rangle =   \langle \rangle f_1(t) , g_1(t) \rangle
\]
giving the position of the planet's center relative to the center of the sun ($O$) in the $xy$-coordinate system in terms of the number of years since the start of the year.

\item Express the components of the vector
 \[
  \overrightarrow{PQ} =  \langle u, v \rangle =   \langle \rangle f_2(t) , g_2(t) \rangle
\]
giving the position of point $Q$ relative to $P$ in the $uv$-coordinate system in terms of the number of years since the start of the year.

\item Use (a) and (b) to express the vector 
\[
  \overrightarrow{OQ} =  \langle x, y \rangle =   \langle \rangle f_3(t) , g_3(t) \rangle
\]
giving the position of point $Q$ relative to $O$ in the $xy$-coordinate system in terms of the number of years since the start of the year.

\item Repeat parts (a)-(c) assuming the planet about its axis in the opposite sense that it revolves about its sun. (Drag the slider to $w=-1$ in Line 5.
\end{enumerate}





\end{document}